\documentclass[11pt,a4paper]{article}
\usepackage[T1]{fontenc}
\usepackage{lmodern,amsmath,amssymb,microtype,booktabs}
\usepackage[margin=25mm]{geometry}
\usepackage[colorlinks=true,urlcolor=blue,linkcolor=blue]{hyperref}
\newcommand{\ga}{\mathfrak a}
\newcommand{\gb}{\mathfrak b}
\title{A uniform transverse limit of the dS$_3$ crossing kernel\\
and its local action on test functions}
\author{Alan Day \and Chengheng\thanks{AI research assistant (OpenAI).}}
\date{Working research note --- 31 October 2026}
\begin{document}\maketitle
\begin{abstract}
Combining the two removable Gamma poles before estimation gives a
bound uniform through the collision surface. The physical scalar
kernel has a phase-shifted sine approximation with uniform $O(M^{-1})$
error on a fixed transverse interval. It converges to a unit delta
distribution there with respect to Lebesgue transverse measure.
This local statement does not yet control the physical mass measures,
creation amplitudes or noncompact spectral tails of a Bose observable.
\end{abstract}

\section{Scope and the local theorem}
We retain $\ell=1$, $G_N=0$, proper $SO^+(3,1)$, BD scalar observers,
KG matter, a fixed positive observer-label lower bound, and half-period
regulated creation blocks. The scalar matrix element $U_{00}$ keeps
its physical CG-period phase; it is not a separately unitary scalar
transform. The label $M$ has static energy $\sqrt{1+M^2}$.

Use $c=M$, $b=M+u$, $a=M+u+v$, $d=M+v+r$, $p,q>0$, and
\[
 A_j=\tfrac12+i\alpha_j,\qquad
 \alpha=\tfrac12(v+q,p-u,-u-p,v-q),\qquad z=-ir/2 .
\]
Choose $R>0$ with $\max_j|\alpha_j|+R/2\le1/2$.
All assertions are uniform for $|r|\le R$ and compact transfer
families satisfying these conditions, once $M$ is sufficiently large.
Let $L=\log M$ and
\[
 \delta(r)=\Im\left[2\log\Gamma(1-ir/2)
       +\sum_j\{\log\Gamma(A_j+ir/2)-\log\Gamma(A_j)\}\right],
\]
continuously from $\delta(0)=0$. Define
\begin{equation}
 K_M(r)=\frac{\sin(Lr-\delta(r))}{\pi r},\quad r\ne0,
 \qquad K_M(0)=\frac{L-\delta'(0)}{\pi}.
 \label{K}
\end{equation}
We prove the model statement
\begin{equation}
 \sup_{|r|\le R}|U_{00}(M,r)-K_M(r)|=O(M^{-1}).
 \label{uniform}
\end{equation}
Consequently, for $f\in C_c^2(\mathbb R)$ supported in $[-R,R]$,
\begin{equation}
 \int_{-R}^{R}f(r)U_{00}(M,r)\,dr
 =f(0)+O_f(1/\log M)+O(\|f\|_1/M).
 \label{local}
\end{equation}
The coefficient one is with respect to $dr$. No mass Jacobian or
physical spectral weight has been incorporated into \eqref{local}.

\newpage
\section{Merging before bounding}
The preceding LPRS/Wilson decomposition \cite{lprs,wilson} is
\[
 X=c_AF+c_BG,\quad Y=y_AF+y_BH,\quad U_{00}=2\Re(C_+XY).
\]
Set $\ga=zy_A$, $\gb=-zy_B$ and $d_0=M+v$. The Gamma poles disappear:
\begin{align*}
 \ga&=\Gamma(1+z)\frac{\Gamma(1-ib)}{\Gamma(1-ib-z)}
       \frac{\Gamma(A_3-z)\Gamma(A_4-z)}{\Gamma(A_3)\Gamma(A_4)},\\
 \gb&=\Gamma(1-z)\frac{\Gamma(1+i(d_0+r))}{\Gamma(1+id_0-z)}
       \frac{\Gamma(A_1)\Gamma(A_2)}{\Gamma(A_1-z)\Gamma(A_2-z)}.
\end{align*}
Both equal one at $r=0$. Also $F(0)=H(0)$, and
\begin{equation}
 Y=\frac{\ga F-\gb H}{z},\qquad S=\frac{\ga-\gb}{z}.
 \label{merged}
\end{equation}
Both expressions take their continuous derivative values at zero.

Here are explicit real constants, not sampled maxima. Put
$m=M-\max(|u|,|v|)-R>0$, $B=\max(|u|,|v|)+R$, and
\begin{align*}
 C_0&=\cosh(\pi/2)\sqrt{1+R/(2m)}\,
       e^{\pi R/4}(1-e^{-2\pi m})^{-1/2},\\
 P_1(t)&=\gamma+1+\operatorname{arsinh}t,\qquad
 P_{1/2}=\gamma+2\log2+2+\operatorname{arsinh}1,\\
 Q&=P_1(R/2)+2P_1(M+B)+2P_{1/2}.
\end{align*}
Gamma moduli give $|\ga|,|\gb|\le C_0$.
The digamma difference series \cite{digamma}, compared with its
decreasing integral, gives for $\sigma>0$
\[
 |\psi(\sigma+it)|\le|\psi(\sigma)|+\sigma^{-1}
                         +\operatorname{arsinh}(|t|/\sigma).
\]
For example the extra sum is bounded by its first term plus
$\int_\sigma^\infty |t|/[x\sqrt{x^2+t^2}]\,dx$.
Differentiating the displayed regular coefficients therefore gives
$|\ga'|+|\gb'|\le C_0Q$ and $|S|\le2C_0Q$.

Write $F=\sum f_n$, $H=\sum h_n$. On the whole interval,
\[
 |f_n|,|h_n|\le\frac{1}{(n+1)^2}
              \prod_{j=1}^{n}\frac{j^2}{j^2+m^2},
 \qquad |f_n'|\le6H_n|f_n|,\quad |h_n'|\le H_n|h_n|.
\]
Here $H_n=\sum_{j=1}^n j^{-1}$. The derivative bounds follow by
logarithmic differentiation of the Pochhammer products; the numerator
real parts are $1/2$ and denominator real parts are one.
The sequences $H_n/(n+1)^2$ are summable, justifying termwise
differentiation. Since $\ga(f_n)-\gb(h_n)$ vanishes at zero,
the fundamental theorem of calculus removes the apparent $1/r$.
Using $\sum_{n\ge1}H_n/(n+1)^2\le2$ gives
\begin{equation}
 |Y-S|\le R_Y=
 \frac{2C_0\{(\zeta(2)-1)Q+14\}}{1+m^2}
 =O(\log M/M^2).
 \label{RY}
\end{equation}
For the merged sum through $n=N-1$, its exact-arithmetic tail is
\begin{equation}
 T_Y(N)=\frac{2C_0}{N}
 \prod_{j=1}^{N}\frac{j^2}{j^2+m^2}
 \{Q+7(2+\log N)\}.
 \label{tail}
\end{equation}

\newpage
\section{Complete-kernel and coefficient errors}
Let $D_0=C_+c_A$, $E_0=C_+c_B$,
$R_F=(\zeta(2)-1)/(1+m^2)$, $Y_*=2C_0Q+R_Y$, and
\[
 D_*=\frac{C_0}{4\pi}
       \sqrt{\frac{\sinh(\pi R/2)}{\pi R/2}}.
\]
The exact Gamma modulus formula implies $|D_0|\le D_*$.
The preceding whole-series bound $|G|\le B_G$ uses a sinh product
and a monic degree-ten ratio polynomial; it is uniform on the
present compact window, with
$|E_0|B_G=O(M^8e^{-\pi M/2})$.
We obtain the finite-$M$ estimate
\begin{equation}
 |U_{00}-2\Re(D_0S)|
 \le2D_*\{R_Y+R_FY_*\}+2|E_0|B_GY_*.
 \label{fullbound}
\end{equation}
Thus the complete series error is uniformly $O(\log M/M^2)$
plus an exponential term, including $r=0$. Bound \eqref{fullbound}
is for the Gamma-exact approximation, not directly for \eqref{K}.

It remains to make the Gamma expansion uniform. Define
\[
 \widetilde D=D_0M^{-z},\quad
 \widetilde a=\ga M^{-z},\quad \widetilde b=\gb M^{z},
\]
\[
 a_\infty=\Gamma(1+z)e^{-i\pi z/2}
   \frac{\Gamma(A_3-z)\Gamma(A_4-z)}{\Gamma(A_3)\Gamma(A_4)},
\]
\[
 b_\infty=\Gamma(1-z)e^{-i\pi z/2}
   \frac{\Gamma(A_1)\Gamma(A_2)}{\Gamma(A_1-z)\Gamma(A_2-z)}.
\]
The previous physical phase calculation gives $D_\infty$ and
\[
 \widetilde D=D_\infty\left[
 1+\frac{r}{4M}-\frac{ir(3r+4v)}{8M}+O_{C^1}(M^{-2})\right].
\]
Its modulus follows directly from
\[
 \left|\frac{\widetilde D}{D_\infty}\right|^2
 =\frac{b}{b-r/2}
   \frac{1-e^{-2\pi b}}{1-e^{-2\pi(b-r/2)}}.
\]
For the phase and the two other coefficients, shifted Stirling and
the digamma expansion \cite{stirling} give, uniformly in real $r$,
\[
 \|\widetilde D-D_\infty\|_{C^1}
 +\|\widetilde a-a_\infty\|_{C^1}
 +\|\widetilde b-b_\infty\|_{C^1}=O(M^{-1}).
\]
The derivative estimate uses the digamma remainder on the same
sector; it does not follow merely by formally differentiating an
unspecified pointwise big-$O$ term. All small Gamma arguments stay
away from poles, and large arguments stay away from the negative
real cut.

Put $E_A=\widetilde D\widetilde a-D_\infty a_\infty$ and
$E_B=\widetilde D\widetilde b-D_\infty b_\infty$.
At zero both equal
\[
 \epsilon_0=D_0(M,0)-(4\pi)^{-1}=O(M^{-2}),
\]
by the preceding peak-phase identity, while
$|E_A(r)-\epsilon_0|+|E_B(r)-\epsilon_0|=O(|r|/M)$.
Since $|e^{-irL}-1|\le |r|L$,
\[
 \left|\frac{E_A(r)e^{-irL}-E_B(r)}{z}\right|
       =O(M^{-1}+L/M^2)
\]
uniformly through zero. The limiting coefficients reproduce
$K_M$ by the previously proved phase and purely imaginary
nonoscillating branch. Together with \eqref{fullbound}, this proves
\eqref{uniform}.

\newpage
\section{Local distributional action and its precise boundary}
We give the short Fourier argument to fix the limiting normalization.
For real or complex $f\in C_c^2(\mathbb R)$ supported in $[-R,R]$, set
\[
 g(r)=f(r)\cos\delta(r),\qquad
 q(r)=\frac{g(r)-g(0)}r,\qquad
 h(r)=f(r)\frac{\sin\delta(r)}r .
\]
The removable values at zero make $q,h$ continuously differentiable;
$\delta$ is smooth. Then
\[
 \int_{-R}^{R} fK_M\,dr
 =\frac{2f(0)}{\pi}\operatorname{Si}(LR)
  +\frac1\pi\int_{-R}^{R}\{q(r)\sin Lr-h(r)\cos Lr\}\,dr .
\]
Integration by parts and the Dirichlet integral give
\begin{equation}
 \left|\int fK_M\,dr-f(0)\right|
 \le \frac{6|f(0)|/R+\|q'\|_1+\|h'\|_1}{\pi L}.
 \label{testbound}
\end{equation}
For example
$|\int_x^\infty\sin t/t\,dt|\le2/x$ supplies the sine-integral
tail; the endpoint values of $q$ contribute $2|f(0)|/R$,
and $h$ vanishes at both endpoints.
This is the same standard removable-pole/Dirichlet mechanism used
in the Wilson analysis \cite{wilson}; the physical kernel and its
uniform error have been controlled separately here.

Adding \eqref{uniform} proves \eqref{local}. The constant in the
$O(1/M)$ kernel estimate depends on the compact transfer family.
We have not produced an optimized numerical value for that constant.
The explicit estimates \eqref{RY}--\eqref{fullbound} remain available
for the series part, with no arithmetic-rounding certification.

There is also a useful conditional extension. If $f_M$ are all
supported in this fixed interval, their $C^2$ norms are uniformly
bounded, and $f_M(0)\to f_*(0)$, then
\[
 \int f_M(r)U_{00}(M,r)\,dr\longrightarrow f_*(0).
\]
Indeed the quantities in \eqref{testbound} and $\|f_M\|_1$ are then
uniformly bounded. This states what would suffice for one local
piece of a physical wave packet; it does not assert that the actual
creation amplitudes satisfy these hypotheses.

The conclusions do not include integration over the other transfer
variables, the observer spectral Jacobians, a moving lower bound,
noncompact mass tails, or growing-time factors.
The full Bose free-observer limit therefore remains open.
In particular a sharp moving cutoff or a test family with uncontrolled
derivatives cannot be substituted into this theorem without further
estimates. Full Fock domains, finite-$G_N$ dynamics and a microscopic
de Sitter holographic dictionary remain later objectives.

\newpage
\section{Numerical evidence and reproducibility}
The interval tests use $(u,v,p,q)=(0.1,-0.1,0.4,0.4)$, $R=0.25$,
$M=16,64,256$, and
$r=-0.2,-10^{-12},-10^{-24},0,10^{-24},10^{-12},0.2$.
All 21 complete kernels, evaluated at 100 digits, satisfy
\eqref{fullbound} after separately allowing for saved truncation tails.
The merger uses \texttt{expm1} on the logarithmic ratio for $r\ne0$
and the exact digamma series at zero.

Eleven points were compared with the unmerged representation at
140 digits and doubled cutoff. Their absolute differences range
from about $3.14\,10^{-23}$ at low mass to about $10^{-75}$ at tiny
nonzero offsets; the $M=256,r=0$ change is about $3.71\,10^{-98}$.
The explicit allowance $10^{-65}$ concerns empirical rounding only.
All values, including imaginary residuals, are retained.

For a direct local-integral diagnostic, take
$f_k(r)=r^k(1-(r/R)^2)^3$ on $[-R,R]$, zero outside, $k=0,1,2$.
These functions are $C^2$ and their target values are $(1,0,0)$.
The full kernel gives:
\begin{center}
\begin{tabular}{rrrr}
\toprule
$M$ & $\int f_0 U\,dr$ & $\int f_1 U\,dr$ & $\int f_2 U\,dr$\\
\midrule
64 & $0.4758802344$ & $7.45120\,10^{-5}$ & $0.0030419739$\\
256 & $0.5580727442$ & $4.52997\,10^{-5}$ & $0.0034353921$\\
\bottomrule
\end{tabular}
\end{center}
The 24/48-point Gauss refinements use 144 complete complex kernel
evaluations. The largest refinement change is about $3.11\,10^{-39}$,
well above some propagated series tails (as small as $10^{-191}$).
Neither the tail nor the refinement difference certifies total
quadrature/rounding accuracy. These finite masses are visibly far
from the limiting delta action.

Separately, evaluating only the proved sine approximation for
$L=4,8,16,32,64,128$ gives $f_0$ integrals approximately
$0.46605,0.68713,0.94635,1.00246,0.999651,1.000026$.
These are not full-kernel computations at $M=e^L$.
The distinction is recorded in the data schema.

An independent Decimal/Fraction replay passes 446 checks, including
the elementary constants, exact companion polynomials, all saved
quadrature nodes and input hashes. Numerical checks support the
implementation; the uniform and distributional statements rest on
the arguments above. All prior scientific files and frozen snapshots
are retained. These are model calculations built from established
formulas, not a claim of a new complete holographic dual.

\begin{thebibliography}{9}
\bibitem{lprs} J. Liu et al.,
\href{https://arxiv.org/html/1808.00612v3}{arXiv:1808.00612v3},
Eqs.\ (3.35)--(3.38), the two-series crossing kernel.
\bibitem{wilson} W. Groenevelt,
\href{https://arxiv.org/html/math/0306424v1}{arXiv:math/0306424v1},
Propositions 4.4 and 4.7; removable terms and Dirichlet integrals.
\bibitem{digamma} NIST DLMF,
\href{https://dlmf.nist.gov/5.7}{Section 5.7.6}, digamma difference series.
\bibitem{stirling} NIST DLMF,
\href{https://dlmf.nist.gov/5.11}{Section 5.11}, log-Gamma, digamma
and Gamma-ratio asymptotics. Checked 31 October 2026.
\end{thebibliography}
\end{document}
